\section{Stage 3: from many samples to a distribution over factors}\label{sec:agg}
\sectionbrief
{for each posterior sample $\hat u_s$ a core result: factors $Q^{(s)}_j$ with coefficient vectors $\bs\beta^{(s)}_j$, weights $w^{(s)}_j$, the structure posterior, $\gamma_N$, $c_F$, $\rho$.}
{the single-sample outputs differ from sample to sample. The stage must merge them into one object that says which factors appear, how often, how strongly, and \emph{how many} there are.}
{it produces the cloud $\mathcal W$ of the analytic note and the posterior on the number of factors, which later sets the number of groups.}

\begin{symbols}
\begin{tabular}{@{}L{2.6cm}L{6.0cm}L{4.3cm}@{}}
$\bs\beta^{(s)}_j$ & the vector of monic coefficients of factor $j$ in posterior sample $s$, padded with zeros to the shape $(m\le2,\ dx\le4)$ & output of the core (Sec.~\ref{sec:core})\\
$w^{(s)}_j,\ \pi_s$ & weight of factor $j$ in sample $s$ (Eq.~\eqref{eq:wj}); weight of sample $s$ & core; interpolation\\
$\delta_{\bs\beta}$ & ``Dirac point mass'': one unit of probability placed exactly at the point $\bs\beta$ & measure theory\\
$\mathcal W$ & the cloud: all the points $\bs\beta^{(s)}_j$, each with mass $\pi_sw^{(s)}_j$ & built here\\
$W_j,\ \bar{\bs\beta}_j,\ \Sigma_j$ & total (normalised) mass, mass-weighted mean and covariance matrix of the $j$-th cluster of the cloud & statistics of $\mathcal W$\\
$H,\ k_{\mathrm{eff}}$ & entropy $H=-\sum_jW_j\ln W_j$ and $k_{\mathrm{eff}}=e^{H}$: the ``effective number'' of equally weighted factors & information theory\\
$k$ & the number of distinct factors (the unknown that the DP posterior estimates) & the unknown\\
$\gamma_0$ & concentration parameter of the Dirichlet process: small $\gamma_0$ favours few clusters & model hyperparameter\\
ELBO & evidence lower bound: the quantity a variational method maximises; here used to pick $\gamma_0$ & variational inference\\
\end{tabular}
\end{symbols}

\subsection{The cloud $\mathcal W$}
Each factor of each sample is encoded as the vector $\bs\beta$ of its monic coefficients, zero-padded to the shape $(m\le2,\ dx\le4)$, so that equivalent structures (a $(1,2)$ factor with $c_2=0$ and a $(1,1)$ factor) give the same vector. It carries the mass $\pi_sw^{(s)}_j$. The result is the weighted point cloud
\begin{equation}\label{eq:cloud}
\mathcal W=\sum_{s=1}^{N_s}\pi_s\sum_jw^{(s)}_j\ \delta_{\bs\beta^{(s)}_j}.
\end{equation}
From it come the dominant factors (modes of the cloud) with total mass $W_j$, mean $\bar{\bs\beta}_j$ and covariance $\Sigma_j$, and the effective number of factors $k_{\mathrm{eff}}=\exp(-\sum_jW_j\log W_j)$. Per-sample distributions of $k_{\mathrm{eff}}$, $\rho$, $c_F$ and $\gamma_N$ are kept as \emph{distributions} (mean and quantiles), together with the co-occurrence matrix of factors and the structure posterior pooled over samples.

\begin{explained}{why a cloud and not an average}
\textbf{Step 1.} Suppose the true law has two factors, transport with $a=-1$ and diffusion with $D=0.3$. In eight posterior samples the fitted coefficients are $-1.00,\,-0.99,\,-1.01,\dots$ and $0.29,\,0.31,\dots$, with weights near $0.77$ and $0.23$. The cloud has two tight clusters, with masses $\approx0.77$ and $\approx0.23$.

\medskip\textbf{Step 2.} Averaging the coefficients over samples would be meaningless as soon as one sample returns a different structure (for instance a spurious extra factor). The cloud keeps every outcome with its probability.

\medskip\textbf{Step 3: $k_{\mathrm{eff}}$.} For masses $(0.77,0.23)$, $H=-0.77\ln0.77-0.23\ln0.23\approx0.54$, so $k_{\mathrm{eff}}=e^{0.54}\approx1.72$; a single factor gives $k_{\mathrm{eff}}=1$ and two equal ones $2$. The note then set $k=\lceil k_{\mathrm{eff}}\rceil=2$.

\medskip\textbf{Step 4: why that rule is fragile.} If a noisy sample puts $1.4\%$ of the mass on spurious factors, $k_{\mathrm{eff}}$ rises from $2.00$ to $2.14$ and $\lceil k_{\mathrm{eff}}\rceil=3$: the true number of laws is $2$. A ceiling is discontinuous; this motivates the next subsection.
\end{explained}

\subsection{A posterior on $k$: Dirichlet-process mixture}
Instead of $k=\lceil k_{\mathrm{eff}}\rceil$, the number of factors gets a posterior. A truncated stick-breaking Dirichlet-process mixture ($8$ components, diagonal covariances, variational, scikit-learn's Bayesian Gaussian mixture) is fitted on the cloud, separately for each $\tau$-degree $m$ (the degree is an invariant of a factor, so factors of different $m$ never share a cluster). Four details decide whether it works; each was found by a failure.
\begin{enumerate}[nosep]
\item \textbf{Tracks are the DP clusters.} Matching factors across samples by nearest coefficient vector, without a distance gate, merged different laws that share $m$ (transport and heat are both $m=1$): $k_{\mathrm{eff}}$ came out as $1.00$. A hand-set gate would break the design rule; the DP cluster is the gate.
\item \textbf{Every factor once, plus draws.} Replicating each factor $50$ times to carry its mass created fake evidence and split heat into 6 clusters. Now every observed factor enters once, plus $n$ extra draws proportional to $\pi_sw^{(s)}_j$.
\item \textbf{One isotropic scale.} Standardising each coefficient separately inflated negligible ones (a $\xi^3$ coefficient of $10^{-4}$ in some samples and $0$ in others became a bimodal unit-scale direction).
\item \textbf{A merge step.} With few samples a DP over-splits (Miller--Harrison: the number of clusters of a DP mixture is not a consistent estimator of $k$). Clusters that \emph{never co-occur in a sample} are merged when BIC prefers one Gaussian with \emph{full} covariance to two. Full covariance matters because fitted coefficients are strongly correlated ($\operatorname{corr}(c_0,c_2)=0.97$ for heat).
\end{enumerate}
The DP concentration $\gamma_0\in\{10^{-3},\dots,10\}$ is chosen by the variational lower bound (empirical Bayes). The posterior of $k$ is then a Monte Carlo quantity: draw $s\sim\pi$, a cluster label per factor from the responsibilities, and count distinct labels. A finite-mixture BIC posterior over $k$ is computed alongside.

\subsection{What has been measured}
On transport plus diffusion (noise $10^{-3}$ and $10^{-2}$, $N_s=8$): $P(k=2)=1.0$; $W=(0.768,0.232)$ against the true weights; $k_{\mathrm{eff}}=1.72$; $\bar\gamma_N=0.174$ and $\bar c_F=0.826=1-\bar\gamma_N$ as predicted for $N=2$. On heat: $P(k=1)=1.0$ after the merge step. These runs used a \emph{stand-in posterior}: the exact grid plus independent noise draws with $\pi_s=1/S$, not GP samples.

\begin{caution}[title=Thin margins]
The merge test for heat at noise $10^{-2}$ sat at $7.96$ against a threshold of $8.32$ with $N_s=8$: it is close to its decision boundary. More samples, or the GP's own covariance of the coefficients, would make it robust. The report recommends the finite-mixture BIC posterior over groups as the group-count rule instead of $\lceil k_{\mathrm{eff}}\rceil$.
\end{caution}

\begin{explained}{a reader asks: is this where we group, and what is $k_{\mathrm{eff}}$ for?}
Yes, this stage prepares grouping. It answers \emph{how many} groups to look for ($k$) and gives each one a centre and a spread. $k_{\mathrm{eff}}$ is a first, crude count: the number of factors counted by mass, so that a faint spurious factor adds only a little. Because rounding it up is fragile (Step 4 above), the pipeline uses the DP posterior on $k$. Note what it does \emph{not} do: it does not group data points and does not find where in $(t,x)$ one law ends and another begins. Here we group \emph{factor estimates} in coefficient space; the grouping of the domain, by local evidence on rectangles, is Section~\ref{sec:group}.
\end{explained}

\begin{takeaway}
Stage 3 gives the factor distribution $\mathcal W$ and a posterior on how many factors there are. The parameters of the next stage (number of groups, initial centres, variances) are \emph{read from this posterior}, not set by hand.
\end{takeaway}
